\documentclass[12pt,a4paper]{article}

\usepackage[T2A]{fontenc}
\usepackage[utf8]{inputenc}
\usepackage[russian]{babel}
\usepackage{amsmath,amssymb,amsthm,mathtools}
\usepackage{geometry}
\usepackage{microtype}
\usepackage{enumitem}
\usepackage{xcolor}
\usepackage{tikz}
\usepackage{hyperref}

\geometry{left=25mm,right=20mm,top=22mm,bottom=24mm}
\hypersetup{
  colorlinks=true,
  linkcolor=blue!45!black,
  urlcolor=blue!55!black
}
\setlist{nosep}
\setlength{\parindent}{1.25em}
\setlength{\parskip}{0.15em}

\newtheorem{theorem}{Теорема}
\newtheorem{proposition}{Утверждение}
\theoremstyle{definition}
\newtheorem{definition}{Определение}
\theoremstyle{remark}
\newtheorem*{remark}{Замечание}

\DeclareMathOperator{\Sp}{Sp}
\DeclareMathOperator{\diag}{diag}
\DeclareMathOperator{\cond}{cond}
\newcommand{\CC}{\mathbb{C}}
\newcommand{\RR}{\mathbb{R}}
\newcommand{\ZZ}{\mathbb{Z}}
\newcommand{\eps}{\varepsilon}
\newcommand{\norm}[1]{\left\lVert #1\right\rVert}

\title{Свёртки на группе $\ZZ_n$}
\author{Конспект лекции}
\date{12 сентября}

\begin{document}

\maketitle

\section{Группа вычетов и рабочее пространство}

Группа вычетов по модулю $n$ имеет вид
\[
  \ZZ_n=\{0,1,\dots,n-1\}.
\]
Сложение и умножение выполняются как обычно, после чего берётся остаток от
деления на $n$. В лекции рассматриваются векторы из $\CC^n$; их координаты
удобно нумеровать с нуля:
\[
  x=(x_0,x_1,\dots,x_{n-1})^{\mathsf T}.
\]
Именно такая нумерация естественна для циклических сдвигов и вычислений по
модулю $n$.

\section{Матрица циклического сдвига}

\begin{definition}
Матрицей циклического сдвига называется матрица
\[
  V=
  \begin{pmatrix}
    0      & 0      & \cdots & 0      & 1\\
    1      & 0      & \cdots & 0      & 0\\
    0      & 1      & \ddots & \vdots & \vdots\\
    \vdots & \ddots & \ddots & 0      & 0\\
    0      & \cdots & 0      & 1      & 0
  \end{pmatrix}\in\CC^{n\times n}.
\]
\end{definition}

При умножении на вектор она циклически сдвигает координаты на одну позицию:
\[
  Vx=(x_{n-1},x_0,x_1,\dots,x_{n-2})^{\mathsf T}.
\]
Аналогично $V^r$ выполняет циклический сдвиг на $r$ позиций. В частности,
\[
  V^n=I.
\]

\section{Спектр циклического сдвига}

Рассмотрим множество корней $n$-й степени из единицы
\[
  E_n=\{\eps_0,\eps_1,\dots,\eps_{n-1}\},
  \qquad
  \eps_k=e^{2\pi i k/n}
  =\cos\frac{2\pi k}{n}+i\sin\frac{2\pi k}{n}.
\]
Эти числа лежат на единичной окружности и образуют вершины правильного
$n$-угольника. Для них выполнено
\[
  \eps_k^n=1,
  \qquad k=0,1,\dots,n-1.
\]

\begin{figure}[ht]
  \centering
  \begin{tikzpicture}[scale=2.15]
    \draw[->,gray!75] (-1.25,0) -- (1.48,0);
    \draw[->,gray!75] (0,-1.20) -- (0,1.34);
    \node[below right] at (1.24,-0.02) {$\operatorname{Re}z$};
    \node[above left] at (-0.02,1.20) {$\operatorname{Im}z$};
    \draw[thick] (0,0) circle (1);
    \draw[blue!45,thick]
      (0:1) -- (45:1) -- (90:1) -- (135:1) --
      (180:1) -- (225:1) -- (270:1) -- (315:1) -- cycle;
    \foreach \angle in {0,45,...,315}
      \fill[blue!65!black] (\angle:1) circle (0.035);
    \draw[->,blue!65!black,thick] (0,0) -- (45:1);
    \draw[blue!65!black] (0.32,0)
      arc[start angle=0,end angle=45,radius=0.32];
    \node[blue!65!black] at (22.5:0.48) {$\frac{2\pi}{8}$};
    \node[above right] at (0:1) {$\eps_0=1$};
    \node[above right] at (45:1) {$\eps_1$};
    \node[above] at (90:1) {$\eps_2$};
    \node[below right] at (315:1) {$\eps_7$};
    \node[below left] at (-0.08,-1.08) {$|z|=1$};
  \end{tikzpicture}
  \caption{Корни восьмой степени из единицы как пример множества $E_n$}
\end{figure}

\begin{theorem}
Спектр матрицы циклического сдвига совпадает с множеством корней $n$-й
степени из единицы:
\[
  \Sp(V)=E_n.
\]
\end{theorem}

\begin{proof}
В соглашениях лекции характеристический многочлен записывается как
\[
  \varphi_V(\lambda)=\det(V-\lambda I).
\]
Вычтем $\lambda$ из каждого диагонального элемента матрицы $V$:
\[
  V-\lambda I=
  \begin{pmatrix}
    -\lambda & 0        & 0        & \cdots & 0        & 1\\
    1        & -\lambda & 0        & \cdots & 0        & 0\\
    0        & 1        & -\lambda & \ddots & \vdots   & \vdots\\
    \vdots   & \ddots   & \ddots   & \ddots & 0        & 0\\
    0        & \cdots   & 0        & 1      & -\lambda & 0\\
    0        & \cdots   & \cdots   & 0      & 1        & -\lambda
  \end{pmatrix}.
\]
В первой строке этой матрицы ненулевыми являются только элементы
$a_{11}=-\lambda$ и $a_{1n}=1$. Поэтому разложение Лапласа по первой строке
содержит ровно два слагаемых:
\[
  \det(V-\lambda I)
  =(-\lambda)(-1)^{1+1}\det M_{11}
   +1\cdot(-1)^{1+n}\det M_{1n},
\]
где $M_{1j}$ --- минор, полученный удалением первой строки и $j$-го
столбца.

Первый минор является нижней треугольной матрицей:
\[
  M_{11}=
  \begin{pmatrix}
    -\lambda & 0        & \cdots & 0\\
    1        & -\lambda & \ddots & \vdots\\
    \vdots   & \ddots   & \ddots & 0\\
    0        & \cdots   & 1      & -\lambda
  \end{pmatrix},
  \qquad
  \det M_{11}=\underbrace{(-\lambda)\cdots(-\lambda)}_{n-1}
  =(-\lambda)^{n-1}.
\]
Второй минор является верхней треугольной матрицей с единицами на главной
диагонали:
\[
  M_{1n}=
  \begin{pmatrix}
    1      & -\lambda & 0        & \cdots & 0\\
    0      & 1        & -\lambda & \ddots & \vdots\\
    \vdots & \ddots   & \ddots   & \ddots & 0\\
    0      & \cdots   & 0        & 1      & -\lambda\\
    0      & \cdots   & \cdots   & 0      & 1
  \end{pmatrix},
  \qquad
  \det M_{1n}=\underbrace{1\cdot1\cdots1}_{n-1}=1.
\]
Подставляя найденные определители миноров, получаем
\begin{align*}
  \varphi_V(\lambda)
  &=(-\lambda)(-\lambda)^{n-1}
    +(-1)^{n+1}\cdot1\\
  &=(-1)^n\lambda^n+(-1)^{n+1}\\
  &=(-1)^n(\lambda^n-1).
\end{align*}
Следовательно,
\[
  \varphi_V(\lambda)=0
  \quad\Longleftrightarrow\quad
  \lambda^n=1
  \quad\Longleftrightarrow\quad
  \lambda\in E_n.
\]
\end{proof}

\section{Собственные векторы и дискретное преобразование Фурье}

Напомним: если $Av=\lambda v$ и $v\ne0$, то $\lambda$ называется
собственным числом, а $v$ --- соответствующим собственным вектором.

\begin{theorem}
Собственному числу $\lambda=\eps_k$ матрицы $V$ соответствует собственный
вектор
\[
  v_k=
  \begin{pmatrix}
    1 & \eps_k^{-1} & \eps_k^{-2} & \cdots &
    \eps_k^{-(n-1)}
  \end{pmatrix}^{\mathsf T}.
\]
\end{theorem}

\begin{proof}
Решим систему
\[
  (V-\eps_k I)x=0.
\]
В матричном виде она записывается следующим образом:
\[
  \begin{pmatrix}
    -\eps_k & 0       & 0       & \cdots & 0       & 1\\
    1       & -\eps_k & 0       & \cdots & 0       & 0\\
    0       & 1       & -\eps_k & \ddots & \vdots  & \vdots\\
    \vdots  & \ddots  & \ddots  & \ddots & 0       & 0\\
    0       & \cdots  & 0       & 1      & -\eps_k & 0\\
    0       & \cdots  & \cdots  & 0      & 1       & -\eps_k
  \end{pmatrix}
  \begin{pmatrix}
    x_0\\ x_1\\ x_2\\ \vdots\\ x_{n-2}\\ x_{n-1}
  \end{pmatrix}
  =
  \begin{pmatrix}
    0\\0\\0\\\vdots\\0\\0
  \end{pmatrix}.
\]

Поскольку $\eps_k\in\Sp(V)$, по Теореме~1
\[
  \det(V-\eps_k I)=0.
\]
Значит, все $n$ строк матрицы $V-\eps_kI$ линейно зависимы. Покажем при
этом, что последние $n-1$ строк линейно независимы. Ограничим их первыми
$n-1$ столбцами. Получится квадратная матрица
\[
  B=
  \begin{pmatrix}
    1      & -\eps_k & 0        & \cdots & 0\\
    0      & 1       & -\eps_k  & \ddots & \vdots\\
    \vdots & \ddots  & \ddots   & \ddots & 0\\
    0      & \cdots  & 0        & 1      & -\eps_k\\
    0      & \cdots  & \cdots   & 0      & 1
  \end{pmatrix}.
\]
Она верхняя треугольная, поэтому
\[
  \det B=1\cdot1\cdots1=1\ne0.
\]
Следовательно, строки $B$, а значит, и соответствующие строки исходной
матрицы, линейно независимы. Это можно увидеть и непосредственно: если их
линейная комбинация с коэффициентами $c_1,\dots,c_{n-1}$ равна нулю, то из
первого столбца следует $c_1=0$, из второго --- $c_2=0$, и так далее, пока
не получим $c_1=\dots=c_{n-1}=0$.

Итак, ранг матрицы не меньше $n-1$. С другой стороны, её определитель равен
нулю, поэтому ранг меньше $n$. Следовательно,
\[
  \operatorname{rank}(V-\eps_kI)=n-1,
\]
а первая строка является линейной комбинацией остальных. При решении системы
её можно временно отбросить.

Положим $x_0=1$. Уравнения, отвечающие строкам со второй по последнюю,
имеют вид
\[
  x_0-\eps_kx_1=0,\qquad
  x_1-\eps_kx_2=0,\qquad\dots,\qquad
  x_{n-2}-\eps_kx_{n-1}=0.
\]
Поскольку $\eps_k\ne0$, отсюда последовательно получаем
\[
  x_1=\eps_k^{-1},\quad x_2=\eps_k^{-2},\quad\dots,\quad
  x_{n-1}=\eps_k^{-(n-1)}.
\]
Осталось проверить отброшенное первое уравнение
$-\eps_kx_0+x_{n-1}=0$. После подстановки найденных координат его левая
часть равна
\[
  -\eps_k+\eps_k^{-(n-1)}
  =\eps_k^{-(n-1)}(1-\eps_k^n)=0,
\]
поскольку $\eps_k^n=1$. Таким образом, найденный ненулевой вектор
действительно является решением всей системы.
\end{proof}

Собственные числа $\eps_0,\dots,\eps_{n-1}$ попарно различны, поэтому
соответствующие собственные векторы линейно независимы и образуют базис
пространства $\CC^n$:
\[
  v_0=(1,1,\dots,1)^{\mathsf T},\qquad
  v_1=(1,\eps_1^{-1},\dots,\eps_1^{-(n-1)})^{\mathsf T},\quad\dots
\]

Матрица, составленная из этих векторов-столбцов,
\[
  F=(v_0\ v_1\ \cdots\ v_{n-1})
   =\begin{pmatrix}
      1 & 1 & \cdots & 1\\
      1 & \eps_1^{-1} & \cdots & \eps_{n-1}^{-1}\\
      \vdots & \vdots & \ddots & \vdots\\
      1 & \eps_1^{-(n-1)} & \cdots & \eps_{n-1}^{-(n-1)}
    \end{pmatrix},
\]
называется матрицей дискретного преобразования Фурье.

Переход к базису из собственных векторов диагонализует $V$:
\[
  F^{-1}VF=\diag(\eps_0,\eps_1,\dots,\eps_{n-1}),
\]
или, что то же самое,
\[
  V=F\diag(\eps_0,\eps_1,\dots,\eps_{n-1})F^{-1}.
\]

\section{Циклические матрицы и их символы}

Циклическая матрица с первым столбцом $(a_0,a_1,\dots,a_{n-1})^{\mathsf T}$
имеет вид
\[
  A=
  \begin{pmatrix}
    a_0     & a_{n-1} & a_{n-2} & \cdots & a_1\\
    a_1     & a_0     & a_{n-1} & \cdots & a_2\\
    a_2     & a_1     & a_0     & \ddots & \vdots\\
    \vdots  & \vdots  & \ddots  & \ddots & a_{n-1}\\
    a_{n-1} & a_{n-2} & \cdots  & a_1    & a_0
  \end{pmatrix}.
\]

\begin{definition}
Многочлен
\[
  a(t)=a_0+a_1t+a_2t^2+\dots+a_{n-1}t^{n-1}
\]
называется \emph{символом} циклической матрицы $A$.
\end{definition}

\begin{theorem}
Для циклической матрицы $A$ с символом $a(t)$ выполнено
\[
  A=a(V)=a_0E+a_1V+a_2V^2+\dots+a_{n-1}V^{n-1},
\]
где $E=I_n$ --- единичная матрица.
\end{theorem}

\begin{proof}
Посмотрим последовательно на степени матрицы сдвига. Нулевая степень есть
единичная матрица:
\[
  V^0=E=
  \begin{pmatrix}
    1&0&0&\cdots&0\\
    0&1&0&\cdots&0\\
    0&0&1&\ddots&\vdots\\
    \vdots&\vdots&\ddots&\ddots&0\\
    0&0&\cdots&0&1
  \end{pmatrix}.
\]
Первая и вторая степени имеют вид
\[
  V=
  \begin{pmatrix}
    0&0&\cdots&0&1\\
    1&0&\cdots&0&0\\
    0&1&\ddots&\vdots&\vdots\\
    \vdots&\ddots&\ddots&0&0\\
    0&\cdots&0&1&0
  \end{pmatrix},
  \qquad
  V^2=
  \begin{pmatrix}
    0&0&\cdots&1&0\\
    0&0&\cdots&0&1\\
    1&0&\cdots&0&0\\
    0&1&\ddots&\vdots&\vdots\\
    \vdots&\ddots&\ddots&0&0
  \end{pmatrix}.
\]
Матрица $V$ сдвигает координаты на одну позицию, $V^2$ --- на две, и в
общем случае $V^r$ --- на $r$ позиций. Последняя различная степень сдвига
равна
\[
  V^{n-1}=V^{-1}=
  \begin{pmatrix}
    0&1&0&\cdots&0\\
    0&0&1&\ddots&\vdots\\
    \vdots&\ddots&\ddots&\ddots&0\\
    0&\cdots&0&0&1\\
    1&0&\cdots&0&0
  \end{pmatrix},
  \qquad V^n=E.
\]

Это можно записать точно через стандартный базис
$e_0,e_1,\dots,e_{n-1}$:
\[
  V^re_j=e_{(j+r)\bmod n}.
\]
Следовательно, при нумерации строк и столбцов от $0$ до $n-1$
\[
  (V^r)_{ij}=
  \begin{cases}
    1, & i-j\equiv r\pmod n,\\
    0, & i-j\not\equiv r\pmod n.
  \end{cases}
\]
В каждой позиции $(i,j)$ ровно одна из матриц
$E,V,V^2,\dots,V^{n-1}$ содержит единицу: это матрица $V^r$ с
$r=(i-j)\bmod n$. Поэтому элемент суммы
\[
  a_0E+a_1V+a_2V^2+\dots+a_{n-1}V^{n-1}
\]
в позиции $(i,j)$ равен
\[
  a_{(i-j)\bmod n}.
\]
Значит, вся сумма имеет вид
\[
  \begin{pmatrix}
    a_0     & a_{n-1} & a_{n-2} & \cdots & a_1\\
    a_1     & a_0     & a_{n-1} & \cdots & a_2\\
    a_2     & a_1     & a_0     & \ddots & \vdots\\
    \vdots  & \vdots  & \ddots  & \ddots & a_{n-1}\\
    a_{n-1} & a_{n-2} & \cdots  & a_1    & a_0
  \end{pmatrix}=A.
\]
По определению символа правая часть исходного равенства есть именно
$a(V)$, что и требовалось доказать.
\end{proof}

\section{Теорема об отображении спектра}

\begin{theorem}[об отображении спектра]
Пусть $A$ --- линейный оператор в конечномерном комплексном пространстве,
а $p(t)$ --- произвольный многочлен. Тогда
\[
  \Sp\bigl(p(A)\bigr)=p\bigl(\Sp(A)\bigr).
\]
\end{theorem}

\begin{proof}
Докажем два включения.

\smallskip
\noindent
\emph{Первое включение.}
Пусть $\lambda_0\in\Sp(A)$. Тогда существует $v_0\ne0$ такое, что
\[
  Av_0=\lambda_0v_0.
\]
Следовательно, $A^rv_0=\lambda_0^rv_0$ для всех $r\geqslant0$. Если
\[
  p(t)=c_0+c_1t+\dots+c_mt^m,
\]
то после умножения этих равенств на коэффициенты $c_r$ и сложения получаем
\[
  p(A)v_0=p(\lambda_0)v_0.
\]
Значит, $p(\lambda_0)\in\Sp(p(A))$, и потому
\[
  p(\Sp(A))\subseteq\Sp(p(A)).
\]

\smallskip
\noindent
\emph{Обратное включение.}
Пусть $\mu\in\Sp(p(A))$. Тогда матрица $p(A)-\mu I$ необратима. Над полем
$\CC$ многочлен $p(t)-\mu$ раскладывается на линейные множители:
\[
  p(t)-\mu=c_m(t-t_1)(t-t_2)\cdots(t-t_m).
\]
Поэтому
\[
  p(A)-\mu I
  =c_m(A-t_1I)(A-t_2I)\cdots(A-t_mI).
\]
Здесь достаточно рассмотреть неконстантный $p$: для постоянного многочлена
утверждение очевидно.
Если произведение матриц необратимо, то хотя бы один из множителей
$A-t_kI$ необратим. Следовательно, $t_k\in\Sp(A)$. Кроме того,
$p(t_k)-\mu=0$, то есть $\mu=p(t_k)$. Значит,
\[
  \Sp(p(A))\subseteq p(\Sp(A)).
\]
Оба включения доказаны.
\end{proof}

\section{Темы докладов}

\subsection{Нормы матриц}

Необходимо рассмотреть три подчинённые матричные нормы:
\begin{align*}
  \norm{A}_1
    &=\max_{1\leqslant j\leqslant n}\sum_{i=1}^{n}|a_{ij}|,
    &&\text{максимальная сумма модулей по столбцам},\\
  \norm{A}_{\infty}
    &=\max_{1\leqslant i\leqslant n}\sum_{j=1}^{n}|a_{ij}|,
    &&\text{максимальная сумма модулей по строкам},\\
  \norm{A}_2
    &=\sqrt{\lambda_{\max}(A^*A)},
    &&\text{спектральная норма}.
\end{align*}
Требуется доказать оценку
\[
  \boxed{\norm{A}_2\leqslant
  \sqrt{\norm{A}_1\norm{A}_{\infty}}}.
\]

\subsection{Число обусловленности матрицы}

Для обратимой матрицы $A$ число обусловленности относительно выбранной
матричной нормы определяется формулой
\[
  \cond(A)=\norm{A}\,\norm{A^{-1}}.
\]
В докладе нужно объяснить:
\begin{itemize}
  \item как при решении системы $Ax=b$ погрешности $A$ и $b$ влияют на
        относительную погрешность решения $x$;
  \item почему величина $\cond(A)$ служит коэффициентом усиления этих
        погрешностей;
  \item почему матрицы с большим числом обусловленности называются плохо
        обусловленными;
  \item почему проблема обусловленности возникает не только у матриц
        большого размера.
\end{itemize}

\end{document}
